\subsection{The initial wave over a stationary shear}\label{cc:uf:sec:first-wave}
For $0<\rho\le1$ let
\begin{equation}\label{cc:uf:base:eq:stationary}
 U_0(x)=\sin(2\pi x_2)(e_1+\rho e_3),\qquad B_\rho=2\pi(e_1+\rho e_3)\otimes e_2
\end{equation}
be the stationary shear of Section~\ref{cc:sec:problem} and its gradient
at the origin. The field $U_0$ is smooth, odd, divergence free and
supported on the Fourier modes $\pm e_2$. As it depends only on $x_2$
and has no second component, $U_0\cdot\nabla U_0=0$. Its pressure and
residual therefore vanish. Each particle keeps its second coordinate,
so the flow and its deformation are
\begin{equation}\label{cc:uf:base:eq:deformation}
 X_0(t,a)=a+t\sin(2\pi a_2)(e_1+\rho e_3),\qquad
 (D_aX_0)^{\pm1}(t,a)=\Id\pm t\cos(2\pi a_2) B_\rho.
\end{equation}
The inverse formula uses $B_\rho^2=0$. The origin is a fixed particle and
$|B_\rho|\le4\pi$.

Take the first wave of frequency $\lambda$ over $U_0$ on $[0,T]$ with
\begin{equation}\label{cc:uf:base:eq:first-wave}
 k_*=\lambda e_1,\qquad \widetilde\chi\equiv1,
\end{equation}
so it is fully present at time zero. At the origin solve the seed
equation \eqref{cc:uf:step:eq:seed} with initial direction $e_2$ and
scale the solution so that $\max_{[0,T]}|\mathbf a(t,0)|=\lambda^{-3/4}$.
At time zero $\kappa(0,a)=\lambda e_1$ at every label, so the projection
in \eqref{cc:uf:tower:eq:start} leaves this initial amplitude unchanged.
The initial frame at the origin is $(e_1,e_2,e_3)$, in which
\begin{equation}\label{cc:uf:base:eq:shear-couplings}
 e_1\cdot B_\rho e_2=2\pi,\qquad e_3\cdot B_\rho e_2=2\pi\rho .
\end{equation}
The second wave starts at time zero, so its amplification parameter is
the ratio $\eta(0)/\nu(0)=\rho$ of these couplings in the first wave's
frame. The shear that amplifies it is the leading
gradient $\mathbf a(0,0)\otimes\kappa(0,0)$ contributed by the first
wave in \eqref{cc:uf:tower:eq:central-return}.

\begin{lemma}\label{cc:uf:base:lem:first}
If $T|B_\rho|\le1/8$, the first wave satisfies the two conditions
\eqref{cc:uf:tower:eq:wave-origin} for every $\mathsf D\ge2$, and its
strength at the origin at time zero lies between $\lambda^{1/4}/2$ and
$\lambda^{1/4}$.
\end{lemma}
\begin{proof}
At the origin the seed equation reads
$\mathbf b_t=-(\Id-2P_\kappa)B_\rho\mathbf b$. Since $\Id-2P_\kappa$
is a reflection, every nonzero solution satisfies
$\bigl|\partial_t\log|\mathbf b|\bigr|\le|B_\rho|$. Applying this estimate
to $\mathbf b$ and to the seed $\mathbf a$, whose length is proportional
to $\gamma$, gives
\[
 \frac{|\mathbf b(t)|}{\gamma(t)}
 \le e^{2T|B_\rho|}\frac{|\mathbf b(t')|}{\gamma(t')}
 \le 2\frac{|\mathbf b(t')|}{\gamma(t')}\qquad(0\le t'\le t\le T).
\]
The same estimate on the seed gives
$e^{-T|B_\rho|}\lambda^{-3/4}\le|\mathbf a(0,0)|\le\lambda^{-3/4}$.
Multiplying by $|\kappa(0,0)|=\lambda$ gives the initial strength.

For the sign of the coupling, the frame equations \eqref{cc:eq:frame}
give $|\partial_tE_1|,|\partial_tE_2|\le|B_\rho|$. Thus the first coupling
can change by at most $2t|B_\rho|^2$ from its value in
\eqref{cc:uf:base:eq:shear-couplings}, and
\[
 \nu(t,0)=E_1(t)\cdot B_\rho E_2(t)\ge 2\pi-2t|B_\rho|^2\ge\pi .
\]
\end{proof}
